%%%%%%%%%%%%%%%%%%%%%%%%%%%%%%%%%%%%%%%%%
% Developer CV
% LaTeX Template
% Version 1.0 (28/1/19)
%
% This template originates from:
% http://www.LaTeXTemplates.com
%
% Authors:
% Jan Vorisek (jan@vorisek.me)
% Based on a template by Jan Küster (info@jankuester.com)
% Modified for LaTeX Templates by Vel (vel@LaTeXTemplates.com)
%
% License:
% The MIT License (see included LICENSE file)
%
%%%%%%%%%%%%%%%%%%%%%%%%%%%%%%%%%%%%%%%%%

%----------------------------------------------------------------------------------------
%	PACKAGES AND OTHER DOCUMENT CONFIGURATIONS
%----------------------------------------------------------------------------------------

\hyphenpenalty=10000

\documentclass[9pt]{developercv} % Default font size, values from 8-12pt are recommended

%\newcommand{\jobtitle}{IT-Specialist with focus on: \\ AI \& Data Engineering}


\newcommand{\skillschart}{%
	\begin{barchart}{5.5}
		\baritem{JavaEE}{92}
		\baritem{LLMs}{70}
		\baritem{LangGraph}{58}
		\baritem{Vue/React}{70}
		\baritem{Python}{90}
		\baritem{Kubernetes \& Ansible}{80}
		\baritem{SQL}{92}
	\end{barchart}
}

\newcommand{\toolbubbles}{
	\begin{center}
		\bubbles{4/VSCode, 5/Git, 6/Linux, 6/DevOps, 5/ArgoCD, 4/Ansible}
	\end{center}
}





% ############### consider this 

%\newcommand{\jobtitle}{AI Engineer with focus on: \\ Large Language Models \& GenAI Solutions}

%\newcommand{\skillschart}{%
	%	\begin{barchart}{5.5}
		%		\baritem{Python}{95}
		%		\baritem{LLMs \& Prompt Engineering}{90}
		%		\baritem{LangChain \& LangGraph}{85}
		%		\baritem{RAG Systems}{80}
		%		\baritem{Vector Databases}{75}
		%		\baritem{AWS/Cloud Platforms}{80}
		%		\baritem{API Development}{85}
		%	\end{barchart}
	%}
%
%\newcommand{\toolbubbles}{6/Python, 6/LangChain, 5/LangGraph, 5/AWS, 5/Docker, 4/Git, 4/Pinecone}
\newcommand{\jobtitle}{Data Engineer with focus on \\ Data Platforms \& ML}

\newcommand{\skillschart}{%
	\begin{barchart}{5}
		\baritem{Python \& R}{95}
		\baritem{SQL}{95}
		\baritem{Apache Spark}{85}
		\baritem{AWS (Redshift, Glue, Sagemaker)}{90}
		\baritem{Data Modeling}{95}
		\baritem{Pipeline Orchestration}{75}
		\baritem{Kubernetes}{85}
		\baritem{PowerBI \& Qlik Cloud}{85}
	\end{barchart}
}


\newcommand{\toolbubbles}{
	\begin{center}
		\bubbles{6/Python, 6/SQL, 4/Spark, 5/AWS, 4/Airflow, 6/Docker, 4/DBT}
	\end{center}
}

\newcommand{\titlebeko}{Software Architekt}
\newcommand{\textbeko}{
	At Beko my primary customer is the City of Vienna. There I spearheaded the modernization and feature expansion of a mission-critical Document Management System (DMS). It's a dynamic project where I am working toghether with my team, on the development of dynamic specialized information systems (Fachinformationssysteme / FIS) using Angular, leveraging the core DMS as a robust backend platform for city workflows. I leed the containerization and cloud-native transformation of the legacy application for deployment on Red Hat OpenShift. At last, we migrated core identity and access management systems to Microsoft Entra ID, enhancing security compliance and user authorization workflows. \\ 
	\texttt{OpenShift} \slashsep \texttt{ArgoCD} \slashsep \texttt{Java}  \slashsep \texttt{SpringBoot} \slashsep \texttt{Angular} \slashsep \texttt{Python} \slashsep \texttt{Ansible}
}


% ############### consider this 


%\newcommand{\jobtitle}{Data Engineer with focus on: \\ Cloud-Scale Data Platforms}

%\newcommand{\skillschart}{%
%	\begin{barchart}{5.5}
%		\baritem{Python}{95}
%		\baritem{SQL}{95}
%		\baritem{Apache Spark}{85}
%		\baritem{AWS (Redshift, Glue)}{90}
%		\baritem{Data Modeling}{85}
%		\baritem{Pipeline Orchestration}{85}
%		\baritem{Kubernetes}{75}
%	\end{barchart}
%}

%\newcommand{\toolbubbles}{6/Python, 6/SQL, 5/Spark, 5/AWS, 5/Airflow, 5/Docker, 4/DBT}

%\newcommand{\jobtitle}{IT-Specialist with focus on: \\ AI \& Data Engineering}


\newcommand{\skillschart}{%
	\begin{barchart}{5.5}
		\baritem{JavaEE}{92}
		\baritem{LLMs}{70}
		\baritem{LangGraph}{58}
		\baritem{Vue/React}{70}
		\baritem{Python}{90}
		\baritem{Kubernetes \& Ansible}{80}
		\baritem{SQL}{92}
	\end{barchart}
}

\newcommand{\toolbubbles}{4/VSCode, 5/Git, 6/Linux, 6/DevOps, 5/ArgoCD, 4/Ansible}


% ############### consider this 

%\newcommand{\jobtitle}{Platform Engineer with focus on: \\ Kubernetes \& Cloud Infrastructure}
%
%\newcommand{\skillschart}{%
%	\begin{barchart}{5.5}
%		\baritem{Kubernetes}{95}
%		\baritem{AWS}{90}
%		\baritem{Terraform / IaC}{90}
%		\baritem{Docker}{90}
%		\baritem{CI/CD Pipelines}{85}
%		\baritem{GitOps (ArgoCD)}{85}
%		\baritem{Python Scripting}{80}
%	\end{barchart}
%}
%
%\newcommand{\toolbubbles}{6/Kubernetes, 6/AWS, 5/Terraform, 5/Docker, 5/ArgoCD, 5/Git, 4/Ansible}

%\newcommand{\jobtitle}{IT-Specialist with focus on: \\ AI \& Data Engineering}


\newcommand{\skillschart}{%
	\begin{barchart}{5.5}
		\baritem{JavaEE}{92}
		\baritem{LLMs}{70}
		\baritem{LangGraph}{58}
		\baritem{Vue/React}{70}
		\baritem{Python}{90}
		\baritem{Kubernetes \& Ansible}{80}
		\baritem{SQL}{92}
	\end{barchart}
}

\newcommand{\toolbubbles}{4/VSCode, 5/Git, 6/Linux, 6/DevOps, 5/ArgoCD, 4/Ansible}


% ############### consider this 

%\newcommand{\jobtitle}{Solutions Architect with focus on: \\ Enterprise Cloud Systems}
%
%\newcommand{\skillschart}{%
%	\begin{barchart}{5.5}
%		\baritem{Enterprise Architecture}{90}
%		\baritem{AWS / Azure}{85}
%		\baritem{Kubernetes \& OpenShift}{85}
%		\baritem{JavaEE / Spring}{90}
%		\baritem{System Design}{85}
%		\baritem{Cloud Transformation}{85}
%		\baritem{Security \& Compliance}{80}
%	\end{barchart}
%}
%
%\newcommand{\toolbubbles}{5/AWS, 5/Azure, 5/Kubernetes, 5/OpenShift, 4/Terraform, 5/Git, 4/Docker}


%----------------------------------------------------------------------------------------
 
\begin{document}

%----------------------------------------------------------------------------------------
%	TITLE AND CONTACT INFORMATION
%----------------------------------------------------------------------------------------

\begin{minipage}[t]{0.45\textwidth} % 45% of the page width for name
	\vspace{-\baselineskip} % Required for vertically aligning minipages
	
	% If your name is very short, use just one of the lines below
	% If your name is very long, reduce the font size or make the minipage wider and reduce the others proportionately
	\colorbox{black}{{\HUGE\textcolor{white}{\textbf{\MakeUppercase{Dipl.-Ing.}}}}}
	\colorbox{black}{{\HUGE\textcolor{white}{\textbf{\MakeUppercase{Manuel}}}}} % First name
	
	\colorbox{black}{{\HUGE\textcolor{white}{\textbf{\MakeUppercase{Esberger}}}}} % Last name
	
	\vspace{6pt}
	
	{\huge \jobtitle} % Career or current job title
\end{minipage}
\begin{minipage}[t]{0.275\textwidth} % 27.5% of the page width for the first row of icons
	\vspace{-\baselineskip} % Required for vertically aligning minipages
	
	% The first parameter is the FontAwesome icon name, the second is the box size and the third is the text
	% Other icons can be found by referring to fontawesome.pdf (supplied with the template) and using the word after \fa in the command for the icon you want
	\icon{MapMarker}{12}{Vienna}\\
	\icon{Phone}{12}{+436644115694}\\
	\icon{At}{12}{\href{mailto:esberger.manuel@gmx.at}{esberger.manuel@gmx.at}}\\	
\end{minipage}
\begin{minipage}[t]{0.275\textwidth} % 27.5% of the page width for the second row of icons
	\vspace{-\baselineskip} % Required for vertically aligning minipages
	
	% The first parameter is the FontAwesome icon name, the second is the box size and the third is the text
	% Other icons can be found by referring to fontawesome.pdf (supplied with the template) and using the word after \fa in the command for the icon you want
	\icon{Linkedin}{12}{\href{https://www.linkedin.com/in/manuel-esberger-9473298b}{manuel-esberger}}\\
	\icon{Github}{12}{\href{https://github.com/Heholord}{github.com/heholord}}\\
\end{minipage}

\vspace{0.5cm}

%----------------------------------------------------------------------------------------
%	INTRODUCTION, SKILLS AND TECHNOLOGIES
%----------------------------------------------------------------------------------------

\cvsect{Who Am I?}

\begin{minipage}[t]{0.4\textwidth} % 40% of the page width for the introduction text
	\vspace{-\baselineskip} % Required for vertically aligning minipages
	The most influential event of my life was when I got my first notebook. I got it when I was ten years old. Since then I never stopped learning about various elements of the digital world. I just wanted to know how things work and this curiosity rooted a path for me which leads through many industries and studies. I am curious what the future might bring.
\end{minipage}
\hfill % Whitespace between
\begin{minipage}[t]{0.5\textwidth} % 50% of the page for the skills bar chart
	\vspace{-\baselineskip} % Required for vertically aligning minipages
	\skillschart
\end{minipage}

\toolbubbles

%----------------------------------------------------------------------------------------
%	EXPERIENCE
%----------------------------------------------------------------------------------------

\cvsect{Experience}

\begin{entrylist}
	\entry
		{04/2026 -- Current}
		{Software Architekt}
		{BEKO (Dedicated to the City of Vienna)}
		{At Beko my primary customer is the City of Vienna. There I spearheaded the modernization and feature expansion of a mission-critical Document Management System (DMS). It's a dynamic project where I am working toghether with my team, on connecting the service to a personalized in-house AI Chatbot. I leed the containerization and cloud-native transformation of the legacy application for deployment on Red Hat OpenShift. At last, we migrated core identity and access management systems to Microsoft Entra ID, enhancing security compliance and user authorization workflows. \\ 
			\texttt{LLMs} \slashsep \texttt{LangGraph} \slashsep \texttt{OpenShift} \slashsep \texttt{ArgoCD} \slashsep \texttt{JavaEE}  \slashsep \texttt{SpringBoot} \slashsep \texttt{Python}
		}
	\entry
		{06/2020 -- 03/2026}
		{Development Lead and Integration Specialist}
		{MSD}
		{Merck Sharp \& Dohme primarily creates medicine and therapies for the people in need. To become more efficient in doing so I work with them to get ready for the technologies of tomorrow. This transformation enables MSD Austria to become a true data driven company. To do so we have used advanced analytics techniques, to help our business, making the right decisions. \\ 
			\texttt{LLM \& Langchain} \slashsep \texttt{AWS Sagemaker, Glue \& Terraform} \slashsep \texttt{Apache Airflow} \slashsep \texttt{Dataiku}  \slashsep \texttt{React JS} \slashsep \texttt{Kubernetes} \slashsep \texttt{Redshift \& Apache Spark} \slashsep \texttt{Python \& PyTorch} \slashsep \texttt{R \& Shiny}  \slashsep \texttt{PowerAutomate \& PowerBI}
		}
	\entry
		{07/2018 -- 05/2020}
		{Full Stack Developer / Hybrid Cloud Engineer / DevOps}
		{IBM Client Innovation Center}
		{At IBM I have been located in the ministery of the interior, primarily. I learned how to deal with a rather conservative client and how to open them up to the technologies the present supplies. Within the company I employed and moderated a reoccurring learning event called \texttt{WeLearn}, which gave many IBM employees a platform for an open learning experience. \\
			\texttt{Openshift}\slashsep\texttt{JavaEE}\slashsep\texttt{Vue.js}\slashsep\texttt{Kubernetes}\slashsep\texttt{Python}\slashsep\texttt{DB2}\slashsep\texttt{z/OS}
		}
	\entry
		{10/2016 -- 05/2018\\\footnotesize{part time}}
		{Full Stack Developer / System Admin (Linux)}
		{IBQ}
		{My former employer and I worked on a platform called LearnLinked which enables organizations as well as private persons to schedule and manage their learning courses and certificates. Many customers where related to the public sector, hence opening me the doors to this field which I enjoyed at IBM.\\ 
			\texttt{Python \& Django}\slashsep\texttt{Angular}\slashsep\texttt{Couch DB}\slashsep\texttt{Linux}\slashsep\texttt{Docker \& Container management}}
	\entry
		{04/2014 -- 06/2015}
		{Software Developer}
		{APA-IT}
		{The primary customers for which I worked have been Swarovski, APA (Austrian Press Agency) and the WKO (Wirtschaftskammer Österreich). Swarovski required me to deploy their product on-site in Hall next to Innsbruck. Further, I held developer and editor seminars for the product \textit{Gentics CMS}. \\ 
			\texttt{Spring \& Hibernate}\slashsep\texttt{JavaScript}\slashsep\texttt{PHP}\slashsep\texttt{Postgres}}
	\entry
		{06/2013 -- 09/2013}
		{Back-End Developer}
		{Irian Solutions}
		{Before my military service, I developed the backend server for a small project, which was used by front-end interns. This application was called \textit{Monster Math Expedition}. Since the task was quickly done, I developed a prototype each for a NoSQL and SQL database backend. Close to the end I helped out the front-end development team. \\ 
			\texttt{Spring}\slashsep\texttt{Hibernate}\slashsep\texttt{Bigtable}\slashsep\texttt{Typescript \& PhoneGap}}
\end{entrylist}

%----------------------------------------------------------------------------------------
%	EDUCATION
%----------------------------------------------------------------------------------------

\cvsect{Education and Certificates}

\begin{entrylist}
	\entry
		{10/2019 -- 06/2026}
		{Master's Degree}
		{Technical University of Vienna}
		{\textit{Specialization}: Application Architecture and Kubernetes \\ 
		\textit{Thesis}: "Designing Cloud Applications for Compatibility: A Case Study on Kubernetes and OpenShift"
		}
	\entry
		{09/2019}
		{Red Hat Certified Specialist}
		{certified by Red Hat}
		{in OpenShift Application Development on a Red Hat OpenShift Container Platform 3.6}
	\entry
		{10/2015 -- 09/2018}
		{Bachelor of Science Degree}
		{Technical University of Vienna}
		{in Medical Informatics with a GPA (Grading Point Average) of 1.9. Relevant courseworks are Software development \& Project management, Algorithms and data structures, Mathematics \& Statistics, Data modeling \& Database systems, Logic, Knowledgebased systems, Internet security, Artificial Intelligence and many more.}
	\entry
		{2008 -- 2013}
		{High School Graduation}
		{Vienna Institute of Technology (Ger.: Technologisches Gewerbemuseum Wien)}
		{in Computer Science with an graduation GPA of 1.4 (eq. to 3.72 in U.S.) and final year GPA of 1.16 (e.q. 3.9 in U.S.). While in high school I completed extracurricular courses such as the  attainment of the Cisco CCNA R\&S certificate and a robotic course called Botball.  Relevant courseworks are: Programming \& Algorithms, Distributed Systems/Databases, Project \& Quality Management and many more.}
\end{entrylist}

\cvsect{Internships}

\begin{entrylist}
	\entry
	{07/2012 -- 08/2012}
	{Robotic Researcher}
	{Technical University of Vienna (TU Wien)}
	{The project funded by the TU Wien, required my team to craft hardware with basic Lego and items from the "Botball Educations Robotics Program" set. Meanwhile, I used an industrial camera sponsored by Festo and an ontology to build a network between artificial intelligence agents.\\ 
		\texttt{AI Agents}\slashsep\texttt{Protégé ontology}\slashsep\texttt{Lego robots}\slashsep\texttt{Java}}
	\entry
	{08/2011}
	{Web Application Developer}
	{Irian Solutions}
	{Irian Solutions has been a strong foot of the \textit{Apache MyFaces} open-source JSF project. While my second internship I did some minor bugfixes of this project and had my first active encounter with design patterns and architectural desicions.\\ 
		\texttt{MyFaces}\slashsep\texttt{JSF}\slashsep\texttt{JavaEE}\slashsep\texttt{JPA}}
	\entry
	{07/2010}
	{Software Developer}
	{AWST (Eng.: Applied Sicence, Software and Technology)}
	{At AWST I made my first experiences with projects that have tight time constraints. The application we developed provided an interface for scientific measuring applications. Since the quality had to be consistent, I was assigend to write Unit Tests as well as designing basic form GUIs.  \\ 
	\texttt{Java}\slashsep\texttt{JUnit}\slashsep\texttt{Java Swing}\slashsep\texttt{JPA \& Hibernate}}
\end{entrylist}

\cvsect{Leadership, Presentation Skills and Volnteering Work}

\begin{entrylist}
	\entry
	{10/2017 -- 02/2018}
	{Volunteering Work}
	{R\{C\} -- Refugees Code}
	{Volunteering at a program by the organisation \texttt{R\{C\}} in coorperation with the TU-Wien. The project they have launched focused on helping people in need by making them ready for the IT job market.}
	\entry
	{08/2018}
	{Organisator / Facilitator}
	{Desire to Inspilre}
	{While this one week project I organized with  an intercultural Erasmus+ project about leadership in Villach. This project included 26 participants from 9 different countries.}
	\entry
	{04/2014 -- 10/2014}
	{Project Leader}
	{Gentics}
	{For our costumer Swarovski I took over the lead of the project \textit{Crystal Catalog}. The project successfully went life in July 2014.}
	\entry
	{2012}
	{IEEE Presentation}
	{Global Conference on Educational Robotics (GCER)}
	{On an international conference held in Honolulu, Hawaii called GCER I presented my teams work about autonomous robotics developed in the year prior.}
	\entry
	{2011 -- 2013}
	{Project leader}
	{Vienna Institute of Technology (Ger.: Technologisches Gewerbemuseum Wien)}
	{Throughout the last two years of graduation I took the lead in various software projects. Most of them where agile projects.}
\end{entrylist}

%----------------------------------------------------------------------------------------
%	ADDITIONAL INFORMATION
%----------------------------------------------------------------------------------------

\begin{minipage}[t]{0.3\textwidth}
	\vspace{-\baselineskip} % Required for vertically aligning minipages
	
	\cvsect{Languages}
	
	
	\textbf{German} - native\\
	\textbf{English} - fluent\\
\end{minipage}
\hfill
\begin{minipage}[t]{0.6\textwidth}
	\vspace{-\baselineskip} % Required for vertically aligning minipages
	
	\cvsect{Hobbies}
	
	I love to go for long walks, even when it rains. I go on road biking trips and in general like to learn new things ... like my latest hobby, woodworking.
\end{minipage}
\hfill
%\begin{minipage}[t]{0.3\textwidth}
%	\vspace{-\baselineskip} % Required for vertically aligning minipages
%	
%	\cvsect{Non profit}
%	
%	In my spare time I am a board member of an NGO called wEUnite. 
%\end{minipage}

%----------------------------------------------------------------------------------------


\end{document}
